Para a colisão 1577229 (2024), a rede estima 91,8\% de chance de dano à aeronave, cerca de 12 vezes a média da base (7,7\%). Essa previsão só depende das evidências que estão no Markov blanket de \texttt{houve\_dano}: \texttt{fase\_voo = subida}, \texttt{massa\_aeronave = leve}, \texttt{tamanho\_animal = grande} e \texttt{ingestao\_motor = sim}. Conhecidas essas, \texttt{estacao}, \texttt{periodo}, \texttt{altura}, \texttt{ceu}, \texttt{motor}, \texttt{animal}, \texttt{qtd\_atingida} e \texttt{avisado} deixam de fazer diferença. O que mais pesa é \texttt{tamanho\_animal = grande}: sem essa informação, a estimativa cairia para 33,1\%, 58,8 pontos percentuais de diferença. Em seguida vêm \texttt{ingestao\_motor = sim} (sem ela, 62,1\%) e \texttt{massa\_aeronave = leve} (sem ela, 90,0\%). Se \texttt{tamanho\_animal} fosse \texttt{pequeno} em vez de \texttt{grande}, a rede daria 19,8\%, tanto observando a troca quanto impondo-a por intervenção. Se \texttt{ingestao\_motor} fosse \texttt{nao} em vez de \texttt{sim}, a rede daria 62,0\%, tanto observando a troca quanto impondo-a por intervenção. Se \texttt{massa\_aeronave} fosse \texttt{pesada} em vez de \texttt{leve}, a rede daria 65,6\%, tanto observando a troca quanto impondo-a por intervenção.